\documentclass[a4paper]{article}

% Basic packages
% Español (México) con XeLaTeX: fontspec para fuentes Unicode, polyglossia
% para la división silábica y los nombres del documento en la variante mexicana.
\usepackage{fontspec}
\usepackage{polyglossia}
\setdefaultlanguage[variant=mexican]{spanish}
% Los nombres chinos van como «pinyin (original)»: xeCJK da a los caracteres
% Han su propia fuente; sin ella XeLaTeX los omite con un aviso.
% FandolSong no cubre algunos caracteres de nombres (U+73FA); WenQuanYi Zen Hei sí.
\usepackage[AutoFallBack=true]{xeCJK}
\setCJKmainfont{FandolSong-Regular.otf}
\setCJKfallbackfamilyfont{\CJKrmdefault}{WenQuanYi Zen Hei}
% polyglossia no traduce los nombres de listings.
\AtBeginDocument{\providecommand{\lstlistingname}{}\renewcommand{\lstlistingname}{Listado}%
  \providecommand{\lstlistlistingname}{}\renewcommand{\lstlistlistingname}{Índice de listados}}
\usepackage{amsmath, amssymb}
\usepackage{graphicx}
\usepackage[margin=2.5cm]{geometry}
\usepackage[most]{tcolorbox}
\usepackage{etoolbox}
\usepackage{listings}
\usepackage{hyperref}
\usepackage{booktabs}
\usepackage{subcaption}
\usepackage{float}

% Highlight boxes
\newtcolorbox{knowledgebox}[1]{
    enhanced, colback=blue!5!white, colframe=blue!75!black, colbacktitle=blue!75!black,
    coltitle=white, fonttitle=\bfseries, title={#1}, attach boxed title to top left={yshift=-2mm, xshift=2mm},
    boxrule=1pt, sharp corners
}

\newtcolorbox{importantbox}[1]{
    enhanced, colback=yellow!10!white, colframe=yellow!80!black, colbacktitle=yellow!80!black,
    coltitle=black, fonttitle=\bfseries, title={#1}, sharp corners
}

\newtcolorbox{warningbox}[1]{
    enhanced, colback=red!5!white, colframe=red!75!black, colbacktitle=red!75!black,
    coltitle=white, fonttitle=\bfseries, title={#1}, sharp corners
}

% Listings
\lstset{
    language=Python,
    basicstyle=\ttfamily\small,
    keywordstyle=\color{blue},
    stringstyle=\color{red!60!black},
    commentstyle=\color{green!60!black},
    breaklines=true,
    frame=single,
    numbers=left,
    numberstyle=\tiny\color{gray},
    captionpos=b,
    extendedchars=true
}

% Front-page metadata
\newcommand{\notetitle}{Deep Generative Modeling}
\newcommand{\noteauthors}{Elaboradas a partir de la clase de Alexander Amini}
\newcommand{\notedate}{Primavera de 2025}
\newcommand{\videochannel}{Alexander Amini (MIT)}
\newcommand{\videopublishdate}{2025-03-24}
\newcommand{\videoduration}{48 minutos}
\newcommand{\videourl}{https://www.youtube.com/watch?v=SdTZAMDKrNY}
\newcommand{\videocoverpath}{cover.jpg}
\newcommand{\repourl}{https://github.com/hqhq1025/ai-course-notes}

% Footer with repo info
\usepackage{fancyhdr}
\pagestyle{fancy}
\fancyhf{}
\fancyfoot[C]{\thepage}
\fancyfoot[R]{\footnotesize\color{gray}\href{\repourl}{AI Course Notes}}
\renewcommand{\headrulewidth}{0pt}
\renewcommand{\footrulewidth}{0pt}

\begin{document}

\begin{titlepage}
\centering
{\Large Notas del curso\par}
\vspace{1.2cm}
{\huge\bfseries \notetitle\par}
\vspace{0.8cm}
{\large \noteauthors\par}
\vspace{0.3cm}
{\large \notedate\par}
\vspace{1.2cm}

\includegraphics[width=0.82\textwidth,height=0.45\textheight,keepaspectratio]{\videocoverpath}\par

\vfill
\begin{tcolorbox}[width=0.9\textwidth, colback=black!2!white, colframe=black!60, sharp corners]
\textbf{Autor o canal del video}: \videochannel\par
\textbf{Fecha de publicación}: \videopublishdate\par
\textbf{Duración del video}: \videoduration\par
\textbf{Enlace del video}: \href{\videourl}{\nolinkurl{\videourl}}\par
\textbf{Página del proyecto}: \href{\repourl}{\nolinkurl{\repourl}}\par
\end{tcolorbox}
\end{titlepage}

\tableofcontents
\newpage

%% --- Inicio del contenido --- %%

\section{Introducción: la era de la modelación generativa}

Estamos en una época de enorme transformación en la IA generativa (Generative AI). El objetivo central de la modelación generativa no es solo identificar patrones en los datos y tomar decisiones con base en ellos, sino, más importante aún, poder \textbf{generar instancias de datos completamente nuevas a partir de los patrones aprendidos}. Esta idea aparentemente simple ha provocado en los últimos años cambios revolucionarios en numerosos campos, entre ellos la generación de imágenes, la generación de audio y el procesamiento de lenguaje natural.

\begin{figure}[H]
\centering
\includegraphics[width=0.85\textwidth]{slides-images/slide-02.jpg}
\caption{¿Cuál de estos rostros es real? En realidad, los tres son rostros falsos generados por IA\protect\footnotemark}
\end{figure}
\footnotetext{Slides página 2. Fuente: thispersondoesnotexist.com}

La clase comienza con un ejemplo que invita a la reflexión: se muestran tres fotografías de rostros y se le pide a la audiencia que determine cuál es real. Al revelarse la respuesta, resulta que los tres fueron sintetizados por un modelo generativo. Esto ilustra de manera contundente la potencia de los modelos generativos modernos y da pie al tema central de esta sesión: comprender y dominar los fundamentos teóricos de la modelación generativa profunda.

\begin{importantbox}{Objetivo central de la modelación generativa}
Dadas muestras de entrenamiento provenientes de cierta distribución, aprender un modelo $P_{\text{model}}(x)$ capaz de representar esa distribución, de modo que se acerque lo más posible a la distribución real de los datos $P_{\text{data}}(x)$. Los modelos generativos pueden usarse en dos grandes tareas:
\begin{itemize}
    \item \textbf{Estimación de densidad} (Density Estimation): aprender la distribución de probabilidad subyacente de los datos
    \item \textbf{Generación de muestras} (Sample Generation): muestrear a partir de la distribución aprendida para generar datos nuevos
\end{itemize}
\end{importantbox}

\subsection{Aprendizaje supervisado vs. aprendizaje no supervisado}

En el contenido previo del curso nos hemos enfocado principalmente en problemas de \textbf{aprendizaje supervisado} (Supervised Learning): dados datos etiquetados $(x, y)$, el objetivo es aprender la función que mapea la entrada $x$ a la etiqueta $y$. La modelación generativa, en cambio, pertenece al ámbito del \textbf{aprendizaje no supervisado} (Unsupervised Learning): solo contamos con datos $x$, sin etiquetas, y el objetivo es descubrir la estructura subyacente oculta en los datos.

\begin{figure}[H]
\centering
\includegraphics[width=0.85\textwidth]{slides-images/slide-04.jpg}
\caption{Comparación entre aprendizaje supervisado y aprendizaje no supervisado\protect\footnotemark}
\end{figure}
\footnotetext{Slides página 4.}

\begin{table}[H]
\centering
\caption{Comparación entre aprendizaje supervisado y aprendizaje no supervisado}
\begin{tabular}{lll}
\toprule
\textbf{Dimensión} & \textbf{Aprendizaje supervisado} & \textbf{Aprendizaje no supervisado} \\
\midrule
Datos & $(x, y)$, $x$ son los datos, $y$ es la etiqueta & Solo $x$, sin etiquetas \\
Objetivo & Aprender la función de mapeo $x \to y$ & Descubrir la estructura oculta/subyacente de los datos \\
Tareas típicas & Clasificación, regresión, detección de objetos & Agrupamiento, extracción de características, reducción de dimensionalidad \\
\bottomrule
\end{tabular}
\end{table}

\subsection{Principio central de la modelación generativa}

\begin{figure}[H]
\centering
\includegraphics[width=0.85\textwidth]{slides-images/slide-05.jpg}
\caption{Modelación generativa: de la estimación de densidad a la generación de muestras\protect\footnotemark}
\end{figure}
\footnotetext{Slides página 5.}

El problema fundamental de la modelación generativa puede formularse así: ¿cómo aprender una distribución de probabilidad $P_{\text{model}}(x)$ que se acerque lo más posible a la distribución real de los datos $P_{\text{data}}(x)$?

$$P_{\text{model}}(x) \approx P_{\text{data}}(x)$$

Una vez que hemos logrado aprender un modelo de distribución de probabilidad de este tipo, podemos muestrear a partir de él para generar instancias de datos completamente nuevas, nunca antes vistas. Las muestras generadas reflejarán las características de la distribución de los datos de entrenamiento.

\subsection{¿Por qué nos interesan los modelos generativos?}

Los modelos generativos tienen tres grandes escenarios de aplicación en el mundo real:

\textbf{1. Reducción de sesgos (Debiasing)}: los modelos generativos pueden descubrir automáticamente la distribución subyacente de las características en un conjunto de datos. Por ejemplo, ante un conjunto de datos de rostros, el modelo puede detectar si características como el tono de piel, la pose o la iluminación están sobrerrepresentadas o subrepresentadas, lo que ayuda a crear conjuntos de datos más justos y representativos.

\begin{figure}[H]
\centering
\includegraphics[width=0.85\textwidth]{slides-images/slide-06.jpg}
\caption{Uso de modelos generativos para descubrir sesgos en conjuntos de datos\protect\footnotemark}
\end{figure}
\footnotetext{Slides página 6. Fuente: Amini/Soleymani+ AAAI/AIES 2019.}

\textbf{2. Detección de anomalías (Outlier Detection)}: en escenarios de seguridad crítica como la conducción autónoma, los modelos generativos pueden aprender la distribución de probabilidad de los datos para identificar eventos poco frecuentes situados en la cola de la distribución (como condiciones climáticas extremas o peatones), lo que permite detectar estas anomalías y ajustar el comportamiento del sistema.

\begin{figure}[H]
\centering
\includegraphics[width=0.85\textwidth]{slides-images/slide-07.jpg}
\caption{Modelos generativos aplicados a la detección de anomalías en la conducción autónoma\protect\footnotemark}
\end{figure}
\footnotetext{Slides página 7. Fuente: Amini+ IROS 2018.}

\textbf{3. Generación de muestras (Sample Generation)}: esta es la aplicación más conocida: usar modelos generativos para crear instancias de datos completamente nuevas. Esta es la base de aplicaciones de Generative AI como ChatGPT, la generación de imágenes y la generación de video.

\begin{figure}[H]
\centering
\includegraphics[width=0.85\textwidth]{slides-images/slide-08.jpg}
\caption{Diversas aplicaciones impulsadas por modelos generativos: lenguaje natural, imagen y video, biología\protect\footnotemark}
\end{figure}
\footnotetext{Slides página 8.}

\subsection{Resumen de la sección}

En esta sección se introdujeron los conceptos básicos de la modelación generativa y se explicó la transición del aprendizaje supervisado al aprendizaje no supervisado. El objetivo central de los modelos generativos es aprender la distribución de probabilidad de los datos, y pueden usarse en escenarios como la estimación de densidad, la detección de anomalías y la generación de muestras. Esta clase se centrará en dos arquitecturas fundamentales: el \textbf{Variational Autoencoder (VAE)} y la \textbf{Generative Adversarial Network (GAN)}, ambas pertenecientes a los \textbf{Latent Variable Model} (modelos de variables latentes).

\section{Variables latentes y Autoencoder}

\subsection{¿Qué es una variable latente (Latent Variable)?}

Para entender el concepto de variable latente, el curso recurre a la célebre \textbf{alegoría de la caverna} (Allegory of the Cave) de la República de Platón: un grupo de prisioneros está atado dentro de una caverna, de frente a una pared, y lo único que pueden ver son las sombras que la luz del fuego proyecta sobre esa pared a partir de los objetos. Para los prisioneros, la sombra es su realidad, pero la estructura real del objeto —la parte que nunca pueden observar directamente— es la causa fundamental que produce esas sombras.

\begin{figure}[H]
\centering
\includegraphics[width=0.85\textwidth]{slides-images/slide-10.jpg}
\caption{La alegoría de la caverna de Platón como analogía del concepto de variable latente\protect\footnotemark}
\end{figure}
\footnotetext{Slides página 10. Fuente: la República de Platón.}

\begin{knowledgebox}{Analogía filosófica de la variable latente}
En el aprendizaje automático, los datos que observamos son como las sombras en la pared: son la manifestación de ciertos factores subyacentes (variables latentes). Las variables latentes (Latent Variables) son como los objetos reales que proyectan las sombras en la caverna; aunque no se pueden medir directamente, determinan la distribución de los datos que observamos. Uno de los objetivos centrales del modelado generativo es extraer estas variables latentes a partir de los datos observados.
\end{knowledgebox}

\subsection{Principio básico del Autoencoder}

El \textbf{Autoencoder} (autocodificador) es un método de aprendizaje no supervisado que se usa para aprender representaciones de características de baja dimensión a partir de datos de entrenamiento sin etiquetar. Su idea central es: mediante un codificador (Encoder) se comprime el dato de entrada de alta dimensión a un espacio latente (Latent Space) de baja dimensión, y luego un decodificador (Decoder) intenta reconstruir el dato de entrada original a partir del espacio latente.

\begin{figure}[H]
\centering
\includegraphics[width=0.85\textwidth]{slides-images/slide-13.jpg}
\caption{Autoencoder: el codificador mapea el dato $x$ a un espacio latente de baja dimensión $z$\protect\footnotemark}
\end{figure}
\footnotetext{Slides página 13.}

El mapeo que aprende el codificador:

$$\text{Encoder}: x \to z \quad \text{(alta dimensión} \to \text{baja dimensión)}$$

Donde $x$ es el dato de entrada (por ejemplo, una imagen) y $z$ es el vector de variables latentes. La razón de mantener la dimensión de $z$ muy por debajo de la de $x$ es que queremos que la red \textbf{comprima} los datos y extraiga la representación de características más esencial.

\subsection{Aprender el espacio latente mediante la reconstrucción}

Pregunta clave: ¿cómo entrenar al codificador para que aprenda buenas variables latentes $z$ sin contar con etiquetas?

La respuesta es: conectar un \textbf{decodificador} (Decoder) después del codificador y entrenar la red completa para reconstruir el dato de entrada original.

\begin{figure}[H]
\centering
\includegraphics[width=0.85\textwidth]{slides-images/slide-15.jpg}
\caption{Autoencoder completo: codificación, decodificación y reconstrucción\protect\footnotemark}
\end{figure}
\footnotetext{Slides página 15.}

$$\text{Decoder}: z \to \hat{x} \quad \text{(baja dimensión} \to \text{alta dimensión)}$$

El objetivo de entrenamiento es minimizar la distancia entre el dato de entrada $x$ y la salida reconstruida $\hat{x}$:

$$\mathcal{L}(x, \hat{x}) = \|x - \hat{x}\|^2$$

\begin{itemize}
    \item $x$: dato de entrada original
    \item $\hat{x}$: dato reconstruido después de la codificación y la decodificación
    \item $\|\cdot\|^2$: error cuadrático medio (Mean Squared Error)
\end{itemize}

\begin{importantbox}{La naturaleza autosupervisada del Autoencoder}
¡La función de pérdida del Autoencoder \textbf{no necesita ninguna etiqueta externa}! Usa el dato de entrada mismo como señal de aprendizaje y, al minimizar el error de reconstrucción, obliga a la red a aprender una representación comprimida significativa. Esto hace del Autoencoder un método de aprendizaje sin etiquetas \textbf{autosupervisado}.
\end{importantbox}

\subsection{Dimensión del espacio latente y calidad de la reconstrucción}

La dimensión del espacio latente influye directamente en la calidad de la reconstrucción. El autocodificado es, en esencia, una operación de \textbf{compresión}: cuanto más pequeño es el espacio latente, mayor es el cuello de botella de información, y la red se ve obligada a aprender una representación más compacta.

\begin{figure}[H]
\centering
\includegraphics[width=0.85\textwidth]{slides-images/slide-17.jpg}
\caption{Efecto de la dimensión del espacio latente en la calidad de la reconstrucción: 2D vs. 5D vs. dato real\protect\footnotemark}
\end{figure}
\footnotetext{Slides página 17.}

Como se muestra en la figura, cuando el espacio latente es de 2D, el resultado de la reconstrucción es notoriamente borroso; al aumentar a 5D, la calidad de la reconstrucción mejora de manera considerable. Esto refleja el compromiso entre la capacidad de representación y el grado de compresión.

\subsection{Resumen de la sección}

El Autoencoder logra el aprendizaje de características sin etiquetas mediante una arquitectura de codificador-decodificador. Comprime los datos a un espacio latente de baja dimensión y luego reconstruye el dato de entrada original a partir de ese espacio latente. Pero el Autoencoder tradicional presenta un problema clave: el proceso de codificación es \textbf{determinista} (deterministic); dado el mismo dato de entrada, siempre produce la misma variable latente. Esto significa que el espacio latente puede no ser lo bastante suave ni continuo como para usarse eficazmente en la generación de muestras nuevas. Esto da paso a la sección siguiente sobre el Variational Autoencoder.
\section{Variational Autoencoder (VAE)}

\subsection{De la codificación determinista a la probabilística}

La codificación del Autoencoder tradicional es determinista: dado un dato de entrada $x$, el codificador produce un vector latente $z$ fijo. Esto hace que el espacio latente pueda tener "huecos" y regiones discontinuas, de las cuales muestrear produce salidas sin sentido.

\begin{figure}[H]
\centering
\includegraphics[width=0.85\textwidth]{slides-images/slide-20.jpg}
\caption{Flujo de codificación determinista del Autoencoder tradicional\protect\footnotemark}
\end{figure}
\footnotetext{Diapositiva 20 de las slides.}

El \textbf{Variational Autoencoder (VAE)} introduce una mejora probabilística: el codificador ya no produce un vector latente fijo, sino que produce los \textbf{parámetros de una distribución de probabilidad}: un vector de medias $\mu$ y un vector de desviaciones estándar $\sigma$. Después se \textbf{muestrea} de esa distribución para obtener el vector latente $z$.

\begin{figure}[H]
\centering
\includegraphics[width=0.85\textwidth]{slides-images/slide-22.jpg}
\caption{La mejora central del VAE: espacio latente probabilístico\protect\footnotemark}
\end{figure}
\footnotetext{Diapositiva 22 de las slides. Fuente: Kingma+ ICLR 2014.}

\begin{importantbox}{La idea central del VAE}
El VAE es una \textbf{mejora probabilística} del Autoencoder. El codificador produce una media $\mu$ y una desviación estándar $\sigma$, que definen una distribución de probabilidad sobre la variable latente $q_\phi(z|x)$. Al muestrear de esa distribución, el VAE puede generar distintos vectores latentes y así lograr una generación diversa.
\end{importantbox}

\subsection{La función de pérdida del VAE}

El objetivo de optimización del VAE se compone de dos partes:

\begin{figure}[H]
\centering
\includegraphics[width=0.85\textwidth]{slides-images/slide-25.jpg}
\caption{Optimización del VAE: pérdida de reconstrucción + término de regularización\protect\footnotemark}
\end{figure}
\footnotetext{Diapositiva 25 de las slides. Fuente: Kingma+ ICLR 2014.}

$$\mathcal{L}(\phi, \theta, x) = \underbrace{\mathbb{E}_{q_\phi(z|x)}[\log p_\theta(x|z)]}_{\text{Pérdida de reconstrucción (Reconstruction Loss)}} - \underbrace{D_{KL}(q_\phi(z|x) \| p(z))}_{\text{Término de regularización (Regularization Term)}}$$

\begin{itemize}
    \item \textbf{Pérdida de reconstrucción}: mide la calidad con la que el decodificador reconstruye el dato $\hat{x}$ a partir de la variable latente $z$, y favorece que $\hat{x}$ se acerque lo más posible a $x$
    \item \textbf{Término de regularización (KL Divergence)}: mide la distancia entre la distribución de la variable latente $q_\phi(z|x)$ que produce el codificador y la distribución previa $p(z)$, y favorece que el espacio latente tenga una estructura adecuada
    \item $\phi$: los parámetros del codificador
    \item $\theta$: los parámetros del decodificador
\end{itemize}

\subsection{Distribución previa y KL Divergence}

El VAE necesita definir una \textbf{distribución previa} (Prior) $p(z)$ para la variable latente. La opción más común es la distribución normal estándar:

$$p(z) = \mathcal{N}(\mu = 0, \sigma^2 = 1)$$

\begin{figure}[H]
\centering
\includegraphics[width=0.85\textwidth]{slides-images/slide-28.jpg}
\caption{Distribución previa: la distribución normal favorece que la variable latente se distribuya de manera uniforme en el centro del espacio latente\protect\footnotemark}
\end{figure}
\footnotetext{Diapositiva 28 de las slides. Fuente: Kingma+ ICLR 2014.}

\begin{knowledgebox}{¿Por qué elegir una distribución previa normal?}
Elegir $\mathcal{N}(0, 1)$ como distribución previa tiene dos propósitos importantes:
\begin{enumerate}
    \item \textbf{Codificación concentrada}: favorece que el codificador distribuya la variable latente en la región central del espacio latente
    \item \textbf{Evitar "hacer trampa"}: evita que la red agrupe los puntos de datos de distintas categorías en regiones aisladas específicas del espacio latente (es decir, que memorice los datos), y en cambio la obliga a aprender una representación significativa y continua
\end{enumerate}
\end{knowledgebox}

La \textbf{KL Divergence} (divergencia KL) es una medida de la "distancia" entre dos distribuciones de probabilidad. Para una distribución previa normal, la divergencia KL tiene una expresión analítica:

$$D_{KL}(q_\phi(z|x) \| p(z)) = -\frac{1}{2} \sum_{j=1}^{J} \left(1 + \log(\sigma_j^2) - \mu_j^2 - \sigma_j^2\right)$$

donde $J$ es la dimensión de la variable latente, y $\mu_j$ y $\sigma_j$ son, respectivamente, la media y la desviación estándar de la $j$-ésima variable latente.

\subsection{La intuición de la regularización: continuidad e integridad}

La función de la regularización es asegurar que el espacio latente tenga dos propiedades clave:

\begin{figure}[H]
\centering
\includegraphics[width=0.85\textwidth]{slides-images/slide-30.jpg}
\caption{La regularización asegura la continuidad e integridad del espacio latente\protect\footnotemark}
\end{figure}
\footnotetext{Diapositiva 30 de las slides. Fuente: Joseph Rocca.}

\begin{enumerate}
    \item \textbf{Continuidad (Continuity)}: los puntos vecinos en el espacio latente deben producir, al decodificarse, contenido semánticamente similar
    \item \textbf{Integridad (Completeness)}: muestrear desde cualquier posición del espacio latente debe decodificarse en contenido con sentido
\end{enumerate}

\begin{figure}[H]
\centering
\includegraphics[width=0.85\textwidth]{slides-images/slide-32.jpg}
\caption{Tras la regularización, la continuidad y la integridad del espacio latente quedan garantizadas\protect\footnotemark}
\end{figure}
\footnotetext{Diapositiva 32 de las slides. Fuente: Joseph Rocca.}

\begin{warningbox}{El problema del espacio latente sin regularización}
Si no hay regularización (y solo se usa la pérdida de reconstrucción), la red puede aprender un espacio latente con numerosos "huecos", regiones que no están cubiertas por los datos de entrenamiento. Muestrear de esas regiones huecas produce salidas de ruido sin sentido. La regularización, al obligar a que la distribución de la variable latente se acerque a la normal estándar, rellena esos huecos.
\end{warningbox}

\subsection{Reparameterization Trick}

El entrenamiento del VAE enfrenta una dificultad técnica: la variable latente $z$ se obtiene mediante una operación de \textbf{muestreo}, y el muestreo no es diferenciable, así que no se puede retropropagar a través de la capa de muestreo.

\begin{figure}[H]
\centering
\includegraphics[width=0.85\textwidth]{slides-images/slide-35.jpg}
\caption{Reparameterization Trick: hacer diferenciable el muestreo\protect\footnotemark}
\end{figure}
\footnotetext{Diapositiva 35 de las slides. Fuente: Kingma+ ICLR 2014.}

La solución es el \textbf{Reparameterization Trick} (técnica de reparametrización). La idea central es separar la aleatoriedad de la variable latente $z$:

$$z = \mu + \sigma \odot \epsilon, \quad \epsilon \sim \mathcal{N}(0, 1)$$

\begin{itemize}
    \item $\mu$: el vector de medias que produce el codificador (nodo determinista y diferenciable)
    \item $\sigma$: el vector de desviaciones estándar que produce el codificador (nodo determinista y diferenciable)
    \item $\epsilon$: ruido aleatorio muestreado de la distribución normal estándar (una constante aleatoria fija)
    \item $\odot$: multiplicación elemento a elemento
\end{itemize}

\begin{importantbox}{La esencia del Reparameterization Trick}
Al reescribir $z$ como una función determinista de $\mu$ y $\sigma$ más un ruido aleatorio externo $\epsilon$, el gradiente puede retropropagarse con normalidad a través de $\mu$ y $\sigma$. La aleatoriedad queda "externalizada" como una constante muestreada de una distribución fija, y no bloquea el flujo del gradiente. Esto permite entrenar todo el VAE de extremo a extremo mediante retropropagación.
\end{importantbox}

\subsection{Perturbación de la variable latente e interpretabilidad}

Un VAE bien entrenado tiene una propiedad importante: las distintas dimensiones del espacio latente suelen corresponder a \textbf{características interpretables} diferentes.

\begin{figure}[H]
\centering
\includegraphics[width=0.85\textwidth]{slides-images/slide-38.jpg}
\caption{Perturbación de la variable latente: fijar las demás variables y cambiar una sola variable para observar el cambio semántico\protect\footnotemark}
\end{figure}
\footnotetext{Diapositiva 38 de las slides. Fuente: Kingma+ ICLR 2014.}

Por ejemplo, en un VAE entrenado sobre un conjunto de datos de rostros, si se fijan todas las variables latentes y se cambia lentamente el valor de una sola dimensión, puede observarse un cambio suave en algún atributo específico del rostro generado (como la orientación de la cabeza o la expresión). Esto muestra que el VAE aprendió con éxito las características semánticas implícitas en los datos.

\subsection{$\beta$-VAE y el desacoplamiento del espacio latente}

En el espacio latente de un VAE estándar, distintas dimensiones pueden codificar características entrelazadas entre sí. El $\beta$-VAE refuerza el desacoplamiento agregando un hiperparámetro $\beta > 1$ delante del término de regularización:

$$\mathcal{L}(\theta, \phi; x, z, \beta) = \mathbb{E}_{q_\phi(z|x)}[\log p_\theta(x|z)] - \beta \cdot D_{KL}(q_\phi(z|x) \| p(z))$$

\begin{figure}[H]
\centering
\includegraphics[width=0.85\textwidth]{slides-images/slide-40.jpg}
\caption{El $\beta$-VAE logra un mejor desacoplamiento del espacio latente\protect\footnotemark}
\end{figure}
\footnotetext{Diapositiva 40 de las slides. Fuente: Higgins+ ICLR 2017.}

Cuando $\beta = 1$, se reduce al VAE estándar; cuando $\beta > 1$, se restringe con más fuerza el cuello de botella del espacio latente, lo que favorece que cada variable latente codifique de forma independiente una característica semántica distinta (disentanglement), de modo que cambiar un atributo no afecte a los demás.

\subsection{Síntesis del VAE}

\begin{figure}[H]
\centering
\includegraphics[width=0.85\textwidth]{slides-images/slide-45.jpg}
\caption{Síntesis de los puntos centrales del VAE\protect\footnotemark}
\end{figure}
\footnotetext{Diapositiva 45 de las slides.}

Los cinco puntos centrales del VAE:
\begin{enumerate}
    \item \textbf{Representación comprimida}: comprime la representación de alta dimensión del mundo en un espacio latente de baja dimensión que puede aprenderse
    \item \textbf{Aprendizaje no supervisado}: logra el aprendizaje sin etiquetas mediante la pérdida de reconstrucción
    \item \textbf{Reparameterization Trick}: permite entrenar de extremo a extremo redes que contienen una capa de muestreo
    \item \textbf{Variables latentes interpretables}: mediante el análisis de perturbación se entienden las características semánticas correspondientes a cada variable latente
    \item \textbf{Generación de muestras nuevas}: basta con muestrear del espacio latente y decodificar para generar datos nuevos
\end{enumerate}

\subsection{Resumen de la sección}

El VAE es una extensión probabilística del Autoencoder que, al introducir un espacio latente probabilístico y la regularización mediante KL Divergence, resuelve los problemas de discontinuidad e incompletitud del espacio latente del Autoencoder tradicional. El Reparameterization Trick resuelve de manera ingeniosa la dificultad de retropropagar a través de una capa de muestreo. El VAE es el primer modelo generativo profundo presentado, y pertenece a la familia de los Latent Variable Model.

\section{Generative Adversarial Networks (GANs)}

\subsection{De la modelización explícita al muestreo implícito}

El VAE intenta modelar de manera explícita la densidad de probabilidad de los datos $p(x)$. Pero si nuestro objetivo es solo generar muestras nuevas, ¿es realmente necesario modelar la densidad de forma explícita?

\begin{figure}[H]
\centering
\includegraphics[width=0.85\textwidth]{slides-images/slide-48.jpg}
\caption{La idea de GAN: sin modelar la densidad de manera explícita, aprender directamente el mapeo del ruido a los datos\protect\footnotemark}
\end{figure}
\footnotetext{Diapositiva 48.}

\begin{knowledgebox}{La filosofía de diseño de GAN}
La idea central de GAN es: en lugar de modelar $p(x)$ de manera explícita, aprender directamente una transformación de una distribución simple (como el ruido gaussiano) a la distribución compleja de los datos. En concreto, se muestrea de la distribución de ruido simple y luego se entrena una red generadora que transforma el ruido en muestras de datos verosímiles. La pregunta clave se convierte en: ¿cómo entrenar a este generador?
\end{knowledgebox}

\subsection{La idea del entrenamiento adversario de GAN}

GAN propone una solución elegante: introducir dos redes que compiten entre sí: el \textbf{generador} (Generator, $G$) y el \textbf{discriminador} (Discriminator, $D$).

\begin{itemize}
    \item \textbf{Generador $G$}: recibe ruido aleatorio $z$ e intenta generar datos falsos lo más verosímiles posible $x_{\text{fake}} = G(z)$
    \item \textbf{Discriminador $D$}: recibe una muestra de datos (que puede ser real o generada) y produce la probabilidad de que dicha muestra sea real $D(x) \in {[}0, 1{]}$
\end{itemize}

\begin{figure}[H]
\centering
\includegraphics[width=0.85\textwidth]{slides-images/slide-68.jpg}
\caption{Arquitectura de entrenamiento de GAN: la confrontación entre el generador y el discriminador\protect\footnotemark}
\end{figure}
\footnotetext{Diapositiva 68.}

Las dos redes forman un \textbf{juego adversario}:
\begin{itemize}
    \item $G$ intenta generar datos falsos capaces de engañar a $D$
    \item $D$ intenta distinguir correctamente entre los datos reales y los datos falsos generados por $G$
    \item Cuando el entrenamiento alcanza el óptimo global, $G$ es capaz de generar muestras cuya distribución coincide con la de los datos reales
\end{itemize}

\subsection{La intuición del entrenamiento de GAN}

La clase mostró de manera vívida el proceso de entrenamiento de GAN mediante una animación sobre una línea de datos unidimensional:

\begin{figure}[H]
\centering
\includegraphics[width=0.85\textwidth]{slides-images/slide-50.jpg}
\caption{Intuición del entrenamiento de GAN: el generador parte del ruido y se aproxima gradualmente a la distribución de los datos reales\protect\footnotemark}
\end{figure}
\footnotetext{Diapositiva 50.}

\textbf{La iteración del proceso de entrenamiento}:
\begin{enumerate}
    \item Al inicio, el generador produce, a partir de ruido aleatorio, datos falsos de muy baja calidad
    \item El discriminador aprende a distinguir los datos reales (puntos verdes) de los datos falsos (puntos rosas) y asigna una probabilidad alta a los datos reales
    \item Con base en la retroalimentación del discriminador, el generador ajusta la distribución de los datos que produce para que los datos falsos se acerquen más a los reales
    \item El discriminador se actualiza de nuevo para adaptarse a datos falsos mejores
    \item Este proceso se repite hasta que los datos falsos son indistinguibles de los reales
\end{enumerate}

\begin{figure}[H]
\centering
\includegraphics[width=0.85\textwidth]{slides-images/slide-65.jpg}
\caption{Etapa avanzada del entrenamiento: la distribución de los datos falsos (rosa) se acerca a la de los datos reales (verde)\protect\footnotemark}
\end{figure}
\footnotetext{Diapositiva 65.}

\subsection{La función de pérdida de GAN}

El objetivo de entrenamiento de GAN puede expresarse en forma de un juego \textbf{minimax}:

$$\min_G \max_D \; \mathbb{E}_{z,x}\left[\log D(x) + \log(1 - D(G(z)))\right]$$

\begin{figure}[H]
\centering
\includegraphics[width=0.85\textwidth]{slides-images/slide-70.jpg}
\caption{Función de pérdida de GAN: el generador minimiza y el discriminador maximiza el mismo objetivo\protect\footnotemark}
\end{figure}
\footnotetext{Diapositiva 70.}

\begin{importantbox}{Interpretación de la función de pérdida de GAN}
\textbf{Objetivo del discriminador (maximización)}:
\begin{itemize}
    \item $\log D(x)$: para los datos reales $x$, se desea que $D(x) \to 1$ (identificarlos correctamente como reales)
    \item $\log(1 - D(G(z)))$: para los datos generados $G(z)$, se desea que $D(G(z)) \to 0$ (identificarlos correctamente como falsos)
\end{itemize}

\textbf{Objetivo del generador (minimización)}:
\begin{itemize}
    \item Se desea que $D(G(z)) \to 1$, es decir, que el discriminador confunda los datos falsos con reales
    \item Esto equivale a minimizar $\log(1 - D(G(z)))$
\end{itemize}
\end{importantbox}

\begin{warningbox}{La inestabilidad del entrenamiento de GAN}
En la práctica, el entrenamiento adversario de GAN puede ser muy inestable. Las dos redes deben mantener un progreso de entrenamiento relativamente equilibrado: si el discriminador es demasiado fuerte, el generador no recibe una señal de gradiente útil; si el generador es demasiado fuerte, el discriminador no puede ofrecer una retroalimentación efectiva. Entre los problemas comunes están el colapso de modo (Mode Collapse) y las oscilaciones durante el entrenamiento. Trabajos posteriores, como Wasserstein GAN (WGAN), propusieron funciones de pérdida alternativas más estables.
\end{warningbox}

\subsection{GAN como transformador de distribuciones}

Una vez terminado el entrenamiento, se descarta el discriminador y se usa únicamente el generador para producir datos nuevos:

\begin{figure}[H]
\centering
\includegraphics[width=0.85\textwidth]{slides-images/slide-72.jpg}
\caption{GAN ya entrenado: se usa solo el generador para crear datos nuevos a partir del ruido\protect\footnotemark}
\end{figure}
\footnotetext{Diapositiva 72.}

En esencia, el generador de GAN aprende una transformación de distribuciones: un mapeo de la distribución de ruido gaussiano $z \sim \mathcal{N}(0, 1)$ a la distribución de datos objetivo.

\begin{figure}[H]
\centering
\includegraphics[width=0.85\textwidth]{slides-images/slide-75.jpg}
\caption{GAN es un transformador de distribuciones: un mapeo de la distribución de ruido a la distribución de datos\protect\footnotemark}
\end{figure}
\footnotetext{Diapositiva 75. Fuente: Brock+ ICLR 2019.}

Al realizar una interpolación suave en el espacio de ruido, las muestras generadas también transitan suavemente sobre la variedad de datos, lo que indica que el generador aprendió la estructura continua de los datos.

\subsection{Desarrollo y aplicaciones de GAN}

\textbf{Progressive Growing of GANs}: al aumentar gradualmente el número de capas de la red y la resolución, es posible generar imágenes de rostros de alta resolución sumamente verosímiles.

\begin{figure}[H]
\centering
\includegraphics[width=0.85\textwidth]{slides-images/slide-78.jpg}
\caption{Rostros de alta resolución generados por Progressive GAN (todos generados por IA)\protect\footnotemark}
\end{figure}
\footnotetext{Diapositiva 78. Fuente: Karras+ ICLR 2018.}

\textbf{Conditional GAN y Pix2Pix}: al incorporar información condicional en el generador y el discriminador, es posible realizar tareas de traducción de imágenes pareadas, como generar fotografías realistas de escenas urbanas a partir de mapas de segmentación semántica.

\begin{figure}[H]
\centering
\includegraphics[width=0.85\textwidth]{slides-images/slide-80.jpg}
\caption{Conditional GAN (Pix2Pix): traducción de imágenes pareadas\protect\footnotemark}
\end{figure}
\footnotetext{Diapositiva 80. Fuente: Isola+ CVPR 2017.}

\begin{figure}[H]
\centering
\includegraphics[width=0.85\textwidth]{slides-images/slide-82.jpg}
\caption{Resultados de aplicación de Pix2Pix: conversión bidireccional entre mapas e imágenes aéreas\protect\footnotemark}
\end{figure}
\footnotetext{Diapositiva 82. Fuente: Isola+ CVPR 2017.}

\textbf{CycleGAN}: permite realizar conversiones entre dominios sin necesidad de datos pareados. No se limita a las imágenes; también puede aplicarse a ámbitos como la transformación de voz.

\begin{figure}[H]
\centering
\includegraphics[width=0.85\textwidth]{slides-images/slide-85.jpg}
\caption{Aplicación de CycleGAN en la transformación de voz\protect\footnotemark}
\end{figure}
\footnotetext{Diapositiva 85.}

\begin{knowledgebox}{La innovación de CycleGAN}
Un GAN tradicional genera la distribución objetivo a partir de una distribución de ruido. CycleGAN extiende esta idea: en lugar de partir del ruido, parte de una distribución de datos y aprende el mapeo hacia otra distribución de datos. Su innovación clave es el uso de la \textbf{pérdida de consistencia cíclica} (Cycle Consistency Loss), que garantiza que, al convertir del dominio A al dominio B y de regreso al dominio A, se recupere la entrada original. Esto permite entrenar sin necesidad de datos pareados.
\end{knowledgebox}

\subsection{Resumen de la sección}

Mediante el entrenamiento adversario entre el generador y el discriminador, GAN logra el objetivo de generar muestras de alta calidad sin necesidad de modelar la densidad de manera explícita. En esencia, GAN es un transformador de distribuciones que mapea una distribución de ruido simple a una distribución de datos compleja. Aunque el entrenamiento de GAN presenta el reto de la inestabilidad, la calidad de sus muestras generadas (en particular en imágenes) suele superar a la del VAE, y de él se han derivado variantes potentes como Conditional GAN y CycleGAN.

\section{Comparación y relación entre los dos modelos}

\subsection{VAE vs. GAN: comparación de arquitecturas}

\begin{figure}[H]
\centering
\includegraphics[width=0.85\textwidth]{slides-images/slide-46.jpg}
\caption{Dos grandes familias de modelos de variable latente: VAE y GAN\protect\footnotemark}
\end{figure}
\footnotetext{Diapositiva 46.}

\begin{table}[H]
\centering
\caption{Comparación sistemática entre VAE y GAN}
\begin{tabular}{lll}
\toprule
\textbf{Dimensión} & \textbf{VAE} & \textbf{GAN} \\
\midrule
Método de modelado & Modela explícitamente la densidad $p(x)$ & Aprende implícitamente el proceso de muestreo \\
Estructura de red & Codificador + decodificador & Generador + discriminador \\
Método de entrenamiento & Maximiza el ELBO & Juego adversario minimax \\
Espacio latente & Generado por el codificador, estructurado & Espacio de ruido, sin codificador \\
Calidad de generación & Generalmente más borrosa & Generalmente más nítida y definida \\
Estabilidad del entrenamiento & Relativamente estable & Inestable, propenso al colapso de modo \\
Interpretabilidad & Espacio latente interpretable & Espacio latente con interpretabilidad débil \\
\bottomrule
\end{tabular}
\end{table}

\begin{warningbox}{Por qué el VAE genera imágenes borrosas}
El VAE utiliza una pérdida de reconstrucción a nivel de píxel (como el MSE), lo cual hace que el modelo tienda a producir el «promedio» de varios resultados posibles, generando salidas borrosas. El GAN, en cambio, optimiza la calidad de generación mediante la señal adversaria que proporciona el discriminador, sin esta limitación, por lo que suele generar imágenes más nítidas. Sin embargo, al GAN le falta un codificador explícito como el del VAE, lo que dificulta la inferencia y la manipulación del espacio latente.
\end{warningbox}

\subsection{Resumen de la sección}

El VAE y el GAN son dos paradigmas complementarios de modelos generativos. El VAE ofrece un espacio latente estructurado y un entrenamiento estable, pero su calidad de generación es limitada; el GAN ofrece una calidad de generación sobresaliente, pero su entrenamiento es inestable y carece de un espacio latente explícito. En aplicaciones prácticas, se puede elegir el modelo adecuado según las necesidades, o usar arquitecturas híbridas que combinen las ventajas de ambos.

\section{Práctica: Lab 2 — Sistema de detección de rostros}

El curso incluye una práctica (Lab 2) en la que se entrena un VAE con un conjunto de datos de rostros para descubrir automáticamente variables latentes (como el tono de piel, la pose, la iluminación, etc.) y utilizarlas para detectar sesgos y características subrepresentadas en el conjunto de datos.

\begin{figure}[H]
\centering
\includegraphics[width=0.85\textwidth]{slides-images/slide-89.jpg}
\caption{Lab 2: sistema de detección de rostros basado en VAE\protect\footnotemark}
\end{figure}
\footnotetext{Diapositivas, página 89.}

\begin{knowledgebox}{Sentido pedagógico del Lab 2}
Esta práctica combina teoría y práctica:
\begin{itemize}
    \item Construir un modelo VAE para procesar imágenes de rostros
    \item Descubrir características latentes en los datos mediante la perturbación de variables latentes
    \item Detectar sesgos en los datos de entrenamiento (por ejemplo, si ciertos tonos de piel o poses están subestimados)
    \item Utilizar modelos generativos para crear conjuntos de datos más justos y equilibrados
\end{itemize}
El código de la práctica está disponible en \href{https://github.com/MITDeepLearning/introtodeeplearning}{github.com/MITDeepLearning/introtodeeplearning}.
\end{knowledgebox}

\subsection{Resumen de la sección}

La práctica ayuda a los estudiantes a profundizar, mediante la experiencia directa, su comprensión de los VAE, en particular sobre cómo utilizar modelos generativos para descubrir y mitigar problemas de sesgo en los conjuntos de datos, un tema importante en la investigación sobre equidad en la IA.

\section{Síntesis y ampliación}

\subsection{Repaso de los conocimientos clave}

Esta clase presentó de manera sistemática las dos arquitecturas fundamentales del modelado generativo profundo:

\textbf{1. Variational Autoencoder (VAE)}:
\begin{itemize}
    \item Comprime los datos en un espacio latente probabilístico mediante una arquitectura de codificador-decodificador
    \item Función de pérdida = pérdida de reconstrucción + regularización por divergencia KL
    \item El Reparameterization Trick permite un entrenamiento diferenciable de extremo a extremo
    \item El espacio latente tiene continuidad e integridad, lo que permite manipular características de manera interpretable
\end{itemize}

\textbf{2. Generative Adversarial Network (GAN)}:
\begin{itemize}
    \item Juego adversarial minimax entre el generador y el discriminador
    \item No necesita modelar la densidad de forma explícita: aprende directamente la transformación de la distribución
    \item La calidad de generación es excelente, pero la estabilidad del entrenamiento es el reto central
    \item Dio origen a variantes potentes como Progressive GAN, Conditional GAN y CycleGAN
\end{itemize}

\begin{importantbox}{Una perspectiva unificada de los modelos generativos}
Tanto el VAE como el GAN comparten el mismo objetivo: aprender la distribución de probabilidad subyacente de los datos. El VAE logra este objetivo mediante una estimación explícita de la densidad, mientras que el GAN aprende la distribución de los datos de manera implícita a través del entrenamiento adversarial. Ambos métodos tienen ventajas y desventajas propias, y representan dos filosofías fundamentales del modelado generativo: \textbf{modelado explícito} frente a \textbf{muestreo implícito}.
\end{importantbox}

\subsection{Del VAE/GAN a los modelos generativos actuales}

El VAE (2014) y el GAN (2014) presentados en esta clase son dos hitos en el campo de los modelos generativos. Después de ellos, este campo siguió avanzando con rapidez:

\begin{itemize}
    \item \textbf{Diffusion Models} (modelos de difusión, 2020--): generan datos mediante un proceso gradual de eliminación de ruido, y se han convertido en el método dominante para la generación de imágenes (por ejemplo, Stable Diffusion, DALL-E 2/3)
    \item \textbf{Flow-based Models} (modelos de flujo): modelan la densidad de forma explícita mediante transformaciones invertibles, como Glow y RealNVP
    \item \textbf{Autoregressive Models} (modelos autorregresivos): como la serie GPT, que modelan la distribución de probabilidad condicional token por token
    \item \textbf{Score-based Models}: generan muestras estimando el campo de gradiente de la distribución de los datos
\end{itemize}

\begin{knowledgebox}{¿Por qué los Diffusion Models terminaron superando a los demás?}
Los Diffusion Models combinan las ventajas del VAE y del GAN: al igual que el VAE, tienen un proceso de entrenamiento estable (basado en la optimización de la verosimilitud), y al igual que el GAN, son capaces de generar muestras de alta calidad. Mediante un proceso markoviano de eliminación gradual de ruido, los Diffusion Models evitan tanto la inestabilidad de entrenamiento del GAN como el problema de generación borrosa del VAE. Las clases posteriores de MIT 6.S191 explicarán los Diffusion Models con detalle.
\end{knowledgebox}

\subsection{Lecturas adicionales}

\begin{itemize}
    \item Kingma \& Welling, ``Auto-Encoding Variational Bayes,'' ICLR 2014 --- el artículo fundacional del VAE
    \item Goodfellow et al., ``Generative Adversarial Nets,'' NeurIPS 2014 --- el artículo fundacional del GAN
    \item Higgins et al., ``$\beta$-VAE: Learning Basic Visual Concepts with a Constrained Variational Framework,'' ICLR 2017
    \item Isola et al., ``Image-to-Image Translation with Conditional Adversarial Networks (Pix2Pix),'' CVPR 2017
    \item Karras et al., ``Progressive Growing of GANs,'' ICLR 2018
    \item Zhu et al., ``Unpaired Image-to-Image Translation using Cycle-Consistent Adversarial Networks (CycleGAN),'' ICCV 2017
    \item Brock et al., ``Large Scale GAN Training for High Fidelity Natural Image Synthesis (BigGAN),'' ICLR 2019
    \item Ho et al., ``Denoising Diffusion Probabilistic Models,'' NeurIPS 2020 --- Diffusion Models
    \item Página principal del curso MIT 6.S191: \href{https://introtodeeplearning.com}{introtodeeplearning.com}
    \item Código del Lab 2: \href{https://github.com/MITDeepLearning/introtodeeplearning}{github.com/MITDeepLearning/introtodeeplearning}
\end{itemize}

%% --- Fin del contenido --- %%

\end{document}
